%==============================================================================
% Section 8: Evaluation
%==============================================================================
\section{Evaluation}
\label{sec:evaluation}

We evaluate the \vbfc on physical hardware across four dimensions: functional correctness, performance, security, and usability.

\subsection{Experimental Setup}
\begin{itemize}
    \item \textbf{Target Platforms:} ASUS Z690-A (Intel 600 series, \SI{32}{MB} SPI flash), Lenovo ThinkPad X1 Carbon Gen 9 (Intel 500 series, \SI{16}{MB}), Gigabyte B550 AORUS (AMD 500 series, \SI{32}{MB}).
    \item \textbf{VBFC Hardware:} Revision 1.2 PCBs, RP2040 @ \SI{133}{MHz}, W25Q128JV.
    \item \textbf{Tools:} Saleae Logic Pro 16 (100 MS/s) for SPI capture, \texttt{vbfc-host} v0.9.0, \texttt{UEFITool} v0.28.0 for firmware analysis.
    \item \textbf{Baseline:} Stock boot (no interposer), Dediprog SF600 programmer for reference dumps.
\end{itemize}

\subsection{Functional Correctness}

\subsubsection{UEFI Firmware Patching}
We extracted the stock firmware from each platform, modified it to unlock hidden \ac{BIOS} menus (flashing a known \texttt{0x74} $\to$ \texttt{0xEB} patch at the menu suppression conditional), and deployed via \vbfc.

\begin{table}[h]
\centering
\caption{Hidden Menu Unlock Results}
\label{tab:menu_unlock}
\begin{tabular}{@{}llcc@{}}
\toprule
\textbf{Platform} & \textbf{Target Menu} & \textbf{Patch Size} & \textbf{Result} \\
\midrule
ASUS Z690-A     & Advanced $\to$ CPU Power Management & 1 byte & Success \\
ThinkPad X1C9   & Config $\to$ Intel AMT / VT-d & 1 byte & Success \\
Gigabyte B550   & AMD CBS $\to$ CPPC / C-state Control & 2 bytes & Success \\
\bottomrule
\end{tabular}
\end{table}

All three platforms booted successfully with patched firmware. The modified menus were accessible and functional. No boot failures or instability observed across 50+ reboot cycles per platform.

\subsubsection{ME Region Neutralization}
On the ASUS Z690-A, we applied a \SI{1.5}{MB} patch set neutralizing the \ac{ME} firmware (replacing \ac{ME} modules with NOP stubs, preserving \ac{ME} header structure for \ac{PCH} verification). The system booted to OS with \ac{ME} in "AltMeDisable" state (confirmed via \texttt{mei\_info} in Linux). This demonstrates the \vbfc's capacity for large-scale firmware modification.

\subsubsection{Capacity Extension}
We created a \SI{24}{MB} composite firmware image (original \SI{16}{MB} + \SI{8}{MB} custom \ac{UEFI} application volume) and mapped it via shadow map across original + extension flash. The \ac{PCH} successfully read across the boundary at \SI{16}{MB} with no errors, confirming transparent extension beyond the original flash size.

\subsection{Performance Evaluation}

\subsubsection{Boot Time Overhead}
\begin{table}[h]
\centering
\caption{Boot Time Comparison (Power-on to OS Loader)}
\label{tab:boot_time}
\begin{tabular}{@{}lccc@{}}
\toprule
\textbf{Configuration} & \textbf{ASUS Z690-A} & \textbf{ThinkPad X1C9} & \textbf{Gigabyte B550} \\
\midrule
Stock (no interposer) & 12.3 s & 9.8 s & 11.1 s \\
VBFC Bypass (jumper)  & 12.4 s & 9.9 s & 11.2 s \\
VBFC Active (patched) & 12.5 s & 10.0 s & 11.3 s \\
\bottomrule
\end{tabular}
\end{table}

The \vbfc adds \SI{\le 100}{ms} (within measurement noise) to boot time. The arbiter activation (header verification + shadow map load) takes \SI{\approx 1.5}{ms} and overlaps with \ac{PCH} early initialization.

\subsubsection{SPI Transaction Latency}
Using the Saleae analyzer, we measured \ac{SPI} read transaction latency (CS\# fall to first MISO bit) for 4-byte reads:

\begin{figure}[h]
\centering
\begin{tikzpicture}
\begin{axis}[
    width=0.95\columnwidth,
    height=5cm,
    ybar,
    bar width=6pt,
    ylabel={Latency (ns)},
    ymin=0, ymax=120,
    xtick={1,2,3,4},
    xticklabels={Stock, Bypass, Active (Orig), Active (Ext)},
    ymajorgrids=true,
    grid style=dashed,
    legend style={at={(0.5,1.02)},anchor=south,legend columns=2},
    nodes near coords,
    nodes near coords align={vertical},
    every node near coord/.append style={font=\tiny}
]
\addplot[vbfcBlue, fill=vbfcBlue!70] coordinates {(1, 42) (2, 43) (3, 45) (4, 68)};
\addplot[vbfcAccent, fill=vbfcAccent!70] coordinates {(1, 42) (2, 43) (3, 45) (4, 72)};
\legend{ASUS Z690-A, ThinkPad X1C9}
\end{axis}
\end{tikzpicture}
\caption{SPI Read Latency (4-byte Fast Read @ 50 MHz). Active (Orig) = read from original flash through arbiter. Active (Ext) = read from extension flash via QSPI. The \SI{\approx 25}{ns} overhead for extension reads is within one SCK cycle.}
\label{fig:spi_latency}
\end{figure}

The arbiter adds \SIrange{1}{3}{ns} for original flash passthrough (mux propagation) and \SIrange{25}{30}{ns} for extension flash reads (QSPI access), both negligible at \SI{50}{MHz} SCK (\SI{20}{ns} period).

\subsubsection{Throughput}
Sustained read throughput for large sequential reads (e.g., \ac{UEFI} loading \SI{4}{MB} DXE volume):
\begin{itemize}
    \item Original flash passthrough: \SI{25}{MB/s} (limited by \ac{PCH} SCK frequency).
    \item Extension flash (QSPI XIP): \SI{48}{MB/s} (RP2040 QSPI @ \SI{133}{MHz} DDR).
\end{itemize}
The \vbfc does not bottleneck \ac{PCH} reads; the \ac{PCH} SCK rate is the limiting factor.

\subsection{Security Evaluation}

\subsubsection{Image Rejection Tests}
We attempted to boot with: (a) unsigned image, (b) HMAC-corrupted image, (c) rolled-back version, (d) truncated image. All were correctly rejected; the arbiter fell back to original flash and the system booted stock. The fallback latency was \SI{<5}{ms}.

\subsubsection{Side-Channel Resistance}
We performed a preliminary power analysis (ChipWhisperer-Lite, 10k traces) on the HMAC verification routine. The hardware SHA256 accelerator shows constant-time execution (no data-dependent branches). No first-order leakage detected. Higher-order analysis and EM probing are future work.

\subsubsection{Glitch Resistance}
We attempted voltage glitching (Crowbar) during HMAC verification and version check. The \rp brown-out detector (BOD) and watchdog reliably reset the device before arbiter activation. No successful bypass achieved in 5000 attempts.

\subsection{Usability and Developer Experience}
\begin{itemize}
    \item \textbf{Provisioning Time:} \SI{\approx 30}{s} for full image sign + flash + shadow map update (mostly flash write time).
    \item \textbf{Patch Iteration:} \SI{\approx 5}{s} for modify-patch-flash-reboot cycle (staging slot + atomic shadow map switch).
    \item \textbf{Debugging:} Sniffer capture via \texttt{vbfc-host sniff --duration 10 --output trace.csv} provides full transaction log for post-boot analysis.
    \item \textbf{Documentation:} Complete API docs, hardware design files, and usage guides at \url{https://github.com/SowinySoft/Virtual-Bios-Firmware-Controller}.
\end{itemize}